\documentclass[11pt]{article}
\usepackage[margin=25mm]{geometry}
\usepackage{fontspec}
\setmainfont{Latin Modern Roman}
\usepackage{amsmath,amssymb,mathtools}
\usepackage{unicode-math}
\setmathfont{Latin Modern Math}
\usepackage{xurl,hyperref}
\hypersetup{colorlinks=true,urlcolor=blue}
\setcounter{secnumdepth}{0}
\setlength{\parindent}{0pt}
\setlength{\parskip}{5pt}
\setlength{\emergencystretch}{3em}
\newcommand{\sref}[2]{\hyperlink{source:#1:#2}{[S#2]}}
\newcommand{\cataloguescope}[1]{\paragraph{Recorded target.}#1}
\newcommand{\R}{\mathbb R}
\newcommand{\C}{\mathbb C}
\newcommand{\N}{\mathbb N}
\newcommand{\Z}{\mathbb Z}
\newcommand{\Q}{\mathbb Q}
\newcommand{\F}{\mathbb F}
\newcommand{\eps}{\varepsilon}
\begin{document}
\section{Monolithic versus Staggered Magnetomechanical Relaxation}\label{statement}

\subsection{Definitions and mathematical statement}

\paragraph{Shared notation.}$\mathbb N=\{1,2,\ldots\}$; the standard number fields have their usual meanings. All additional conventions are stated locally.


\paragraph{Definitions and conventions.}
Relaxation replaces a nonconvex material energy by an effective energy allowing fine-scale microstructure, commonly through a lower-semicontinuous envelope or an admissible microstructure minimization. Monolithic relaxation treats mechanical and magnetic variables together. Staggered relaxation alternates or successively relaxes subsets of those variables. Equivalence requires equality of the resulting effective energies or solution behavior under stated coupling and admissibility assumptions.

\paragraph{Target retained from the canonical source catalogue.}
\paragraph{Catalogue wording retained for traceability.}
For coupled magnetomechanical systems, should relaxation be formulated monolithically or as a staggered sequence of partial relaxations, and under what conditions are the two procedures equivalent?

\paragraph{Source scope and qualification.}
The ordering, topology and compatibility constraints of the partial relaxations matter. Equality of two numerical minimization runs does not establish equivalence of the variational relaxation procedures.
\subsection{Short English statement}
Decide when combined and successive relaxations give the same effective magnetomechanical model.
\subsection{Research attempt}
\paragraph{Model, date and investigated scope.}
OpenAI GPT-6 Astra (configured model \texttt{gpt-6-astra}), ultra reasoning; generated 11 October 2026. Proved joint versus separate convexification differ for a coupled scalar energy; diagnoses correlated-mixture loss, not the full constrained magnetomechanical relaxation.

\paragraph{Sources and literature check.}
Read the supplied definition, catalogue wording, source qualifications and reference entries for this question. The historical source/context and inherited status are retained; no fresh literature or novelty certification is claimed. No external theorem is used without proof in this attempt; explicit model assumptions are distinguished from facts about the source model.

\paragraph{A separately convex countercheck.} Consider the scalar coupled energy \(W(x,y)=(xy-1)^2\) on \(\mathbb R^2\). For fixed \(y\), its second derivative in \(x\) is \(2y^2\ge0\), and the symmetric statement holds for \(y\). Thus convexifying in one variable at a time does nothing, even after arbitrarily many alternating steps.
The joint convex envelope behaves differently at the origin. Every convex minorant \(F\le W\) obeys
\[
 F(0,0)\le\tfrac12F(1,1)+\tfrac12F(-1,-1)\le0,
\]
whereas the zero function is a convex minorant because \(W\ge0\). The joint convex envelope is therefore zero at the origin, while the original and all staggered separate convexifications equal one there.
\paragraph{Mechanism and limits.} The lowering uses a correlated mixture of the two variables: both change sign together. Convexification in a single variable with the other fixed cannot access those zero-energy states. This gives an exact finite-dimensional failure of a universal monolithic/staggered equivalence.
\paragraph{Source comparison.} Actual magnetomechanical relaxation may involve quasiconvexity, divergence constraints, boundary data and particular weak topologies rather than ordinary convexification. The example is not claimed to be a compatible magnetoelastic microstructure. The first gap is characterizing which correlated mechanical/magnetic mixtures the source admissibility constraints allow and whether separate relaxation realizes them. Equality of computational iterations alone cannot answer that variational question. No necessary-and-sufficient equivalence theorem for the physical system is proved.

\paragraph{Verification and adversarial audit.}
Checked separate second derivatives, both zero-energy endpoints and the exact joint-envelope value at the origin; distinguished convex from constrained/quasiconvex relaxation. This is a GPT-6 Astra Ultra self-audit, not independent human verification. Prompt A drafting followed by prompt B scope/dependency/boundary self-audit. Checked separate second derivatives, both zero-energy endpoints and the exact joint-envelope value at the origin; distinguished convex from constrained/quasiconvex relaxation. The remaining gap is: Characterize admissible joint microstructures and their attainability by the source’s ordered partial relaxations. Unresolved status and zero original-target completion are retained.

Coordinator checked separate convexity and the exact joint-envelope value at the origin; admissibility as a physical microstructure is not assumed. This additional check was model-assisted, not independent human verification.

\paragraph{Tools, process and contributions.}
Python 3 and local file-reading/writing tools for manuscript serialization and structural checks; mathematical calculations were symbolic reasoning, with no numerical experiment or theorem prover. Contribution prompt A was followed by prompt B. GPT-6 Astra with ultra reasoning drafted the research and performed its stated self-audit; any additional coordinator contribution is identified in the verification and audit. A GPT-6-based Codex coordinator assembled the manuscript and checked repository integration; its exact variant and reasoning configuration were not exposed. No human mathematical contribution, independent human proof check or formal proof certification is supplied. Complete derivations and the consulted source references are included here; private reasoning transcripts are not published.

\paragraph{Final status and first remaining gap.}
\textbf{UNRESOLVED}. Characterize admissible joint microstructures and their attainability by the source’s ordered partial relaxations.

\section{Completion Estimate}
\noindent\textbf{COMPLETION: 0\%}
The proved calculation is an elementary auxiliary result or a test of a restricted strategy; no new portion of the full source target is established.

\subsection{Sources}
\begin{itemize}
\item[\hypertarget{source:statement:1}{S1}] ULAM.AI and the UnsolvedMath Contributors, UnsolvedMath, catalogue entry. Catalogue curation under CC BY 4.0; definitions and English exposition added locally.
\url{https://huggingface.co/datasets/ulamai/UnsolvedMath}.
\item[\hypertarget{source:statement:2}{S2}] Mechanics of Materials: Towards Predictive Methods for Kinetics in Plasticity, Fracture, and Damage (2020), Oberwolfach Reports, DOI 10.4171/owr/2020/13, pp. 713--783. Primary report accessed 2026-10-11; relevant PDF pages 27.
\url{https://doi.org/10.4171/owr/2020/13}.
\item[\hypertarget{source:statement:3}{S3}] Rossi--Thomas (2017) (catalogue-listed context; not independently consulted in this preparation).
\url{https://doi.org/10.1137/15M1051567}.
\item[\hypertarget{source:statement:4}{S4}] Crismale--Lazzaroni--Rossi (2023) (catalogue-listed context; not independently consulted in this preparation).
\url{https://doi.org/10.1007/s10231-022-01280-0}.
\end{itemize}
\label{end:statement}
\end{document}
